\section{从 MoE 到 Attention：架构创新为什么重新重要}

本章把 Attention 放进更大的架构史。Transformer 包含 Attention、FFN 和位置编码。FFN 已经被 MoE 改造；位置编码从绝对位置走向相对位置、RoPE 等；Attention 则在 Full、Linear、Sparse、Sliding Window 等路线中继续演化。访谈中的核心判断是：MoE 是近几年最大架构突破，下一块可能就是 Attention。

\begin{figure}[H]
\centering
\includegraphics[width=0.96\textwidth]{moe-to-attention.png}
\caption{从 MoE 到 Attention：FFN 已经被雕成 MoE，Attention 成为下一块可雕区域。自制概念图，依据 00:58:08--01:02:48 对谈内容整理。}
\end{figure}

\begin{knowledgebox}{读图：Transformer 不是铁板一块}
Transformer 是 Attention + FFN 的堆叠。FFN 变成 MoE 后，算力使用方式发生变化；Attention 仍是长上下文瓶颈，因此 Linear/Sparse 等路线试图继续“雕”架构。
\end{knowledgebox}

\subsection{FLOPs、Loss 与架构}

本节把架构选择转成可比较的目标函数。模型公司不是为了新奇而改架构，而是为了在同样计算预算下得到更好的训练和推理结果。

架构创新背后的朴素目标是：给定相同 FLOPs，取得更低 loss。FLOPs 是 floating point operations，即浮点运算量。好的架构不是简单减少计算，而是在相同预算下更有效地表达函数、更好地使用数据、更高效地适配硬件。

\begin{importantbox}{架构创新的目标函数}
在相同数据和算力约束下，架构创新要么降低训练/推理成本，要么提高表达能力，要么改善硬件效率；最理想的是三者同时改善。
\end{importantbox}

\subsection{为什么数据撞墙会让算法创新更重要}

接下来把算法放回产业约束中。只要数据和算力都还能无脑增长，架构创新的重要性就容易被掩盖；当两者都变贵时，算法效率会重新进入中心。

节目开头提到，数据、算法、算力是驱动人工智能的三驾马车。当数据增长变慢、算力受限或昂贵时，算法和架构的价值会重新上升。中国公司算力相对美国更受限，反而更有动力在 MoE、Sparse Attention、Linear Attention、NoPE/RoPE 等方向上做“省算力但不掉性能”的创新。

\begin{knowledgebox}{三驾马车的约束变化}
\begin{tabularx}{\textwidth}{p{0.18\textwidth}p{0.38\textwidth}X}
\toprule
变量 & 过去的主要做法 & 现在的压力 \\
\midrule
数据 & 继续扩大互联网/代码/多模态数据 & 高质量公域数据增速放缓。 \\
算力 & 堆 GPU、扩大训练规模 & 成本、供给和能源约束加重。 \\
算法 & 沿用 Transformer 主干微调 & 需要在相同 FLOPs 下更有效。 \\
\bottomrule
\end{tabularx}
\end{knowledgebox}

\subsection{本章小结}

Attention 之所以重新成为焦点，是因为长上下文和推理成本把全局注意力推到瓶颈位置。MoE 改造了 FFN，下一阶段可能是 Linear/Sparse/Hybrid Attention 改造注意力层。

